% discussion.tex (v10). No subheadings (Sci Rep).
\section{Discussion}\label{sec:discussion}
This study evaluated ALT models in a cell-year design, in which the labels, the climate input and the scores refer to the same cell and year. Three results follow. First, with the climate of the observation year accumulated to the observation day, a gradient-boosting model of the Stefan residual reduced the Stefan error by \ndrefLeblkRoabs{} and \ndrefCablkRoabs{} cm in the Lena Delta and Canada and matched it in Alaska. Second, the same model remained stable when a region was withheld or when only a few labels were available, whereas direct ML failed in the withheld Lena Delta and needed more labels. Third, in the same cells and years the model was more accurate than five published products in the Lena Delta and on par with the two ML products in Alaska and Canada.

The physics-residual structure works because the two parts carry different information. The Stefan term carries the response to the thaw-season climate of each year and extrapolates to new regions and years with a physically bounded form. The residual model corrects the deviation of the regional coefficient that comes from soil and terrain: the gain needed terrain and soil covariates, it was largest in the conditions where the Stefan error was largest, and the attribution of the residual model pointed to terrain in the Lena Delta and soil organic carbon in Canada. Adding Stefan pseudo-labels pulled the learner towards the Stefan model and did not add accuracy, and direct ML, which has to learn the climate response from the labels, lost \ndrefLeregDzabs{} cm of accuracy in the withheld Lena Delta. With all labels of a region, direct ML can reach a similar or lower point error than the residual model; the advantage of the residual structure is that its accuracy does not depend on having many labels from the same region.

The published accuracies of existing products refer to different site sets and validation schemes (for example a random ten-fold RMSE of 21.6 cm for Liu et al.\cite{Liu2024} and a leave-one-site-out RMSE of 60.5 cm for Wei et al.\cite{Wei2026}) and cannot be compared with each other; scoring models and products on the same cells and years removes this ambiguity. The regional differences in this comparison follow the training data of the products. Wei et al. and Liu et al. were trained on CALM, GTN-P and ABoVE observations that lie within a few kilometres of the Alaska label cells, and in Alaska and Canada they were as accurate as our model. The Lena Delta lies about \ndistweiLemed{} km from the nearest training site of Wei et al., and both products overestimate its thin active layers. A model trained without any Lena Delta label was already more accurate there than Wei et al., and with the Lena Delta labels the error fell to \nccLeRoall{} cm, against \nccLeCcifiveall--\nccLeWeiall{} cm for the products. The residual learner also improved a product used as the anchor. Local labels combined with residual learning can therefore be used to adapt either the Stefan model or an existing product to a region in which the product was not trained.

The label experiments give a practical sequence for a new region. Without local labels, the Stefan model with the coefficient of other regions is the safer starting point, because direct ML can fail. With about 25 labels, recalibrating the Stefan coefficient and learning the residual restored the accuracy in the Lena Delta, and in the external CALM regions five labels reduced the error from \nftAllRszero{} to \nftAllRsfive{} cm, below the ESA CCI v5 product. The external sites also show the limit of the transfer: West and Central Siberia and Svalbard have thaw depths beyond the range of the study regions, and the gain there came from recalibrating the coefficient rather than from the residual model. Canada shows the opposite case, in which the source-trained direct ML was more accurate and recalibration on a small Canadian sample increased the error. The number of local labels needed therefore depends on how far the new region lies from the training range, and 25--100 labels covering the range of soils and terrain are a reasonable target for regions similar to the study regions.

The year-labelled maps are the second product of the design. Each map is driven by the ERA5-Land climate of a stated year, so that a map can be compared with observations or other products of the same year and anomalies can be followed from year to year (for example \nanmaxAk{} cm in Alaska in \nanmaxyrAk). The maps report the area of applicability and a 90\% prediction interval, which mark the cells where the prediction leaves the training conditions. The predicted interannual range is smaller than that of the observed cell means (see also the long-record cells in Supplementary Fig.~S10), partly because the observed cells change between campaigns and partly because the residual model is trained mainly on spatial contrasts.

Several limits apply. Half of the Alaska labels come from 2014, so the effect of the year-specific input could not be resolved there. The Lena Delta labels come from one dataset in \nnblkLe{} blocks, which widens the intervals. The climate input is monthly and the soil covariates are static, so short events and changes in soil moisture are not represented. Active layers thicker than about 3 m, mostly in mountains and at the southern margin of the permafrost zone, lie outside the training range. Daily climate inputs, snow and soil-moisture information and labels from a fourth region such as West Siberia, where both CALM and campaign data exist, are the next steps.

In summary, a Stefan model recalibrated with local labels and corrected by a residual learner gives year-specific ALT that is more accurate than the Stefan model where soil and terrain matter, stable in regions without labels and adaptable with a few labels. Evaluated in the same cells and years, it is more accurate than published products in a region they did not learn and on par where they did. The same procedure can be applied to other permafrost regions and to existing products as anchors, and it provides year-labelled 1 km maps with a stated area of applicability for monitoring the active layer.
